% Lección: funciones continuas (castellano).

\begin{frame}
\didactatitle{Funciones continuas}

\begin{definition}[Continuidad]
Sea $f$ una función definida en un intervalo $I$ y sea $a \in I$. Decimos que
$f$ es \keyterm{continua en $a$} si para cada $\varepsilon > 0$ existe un
$\delta > 0$ tal que
\[
  x \in I, \ |x - a| < \delta \implies |f(x) - f(a)| < \varepsilon .
\]
Es \keyterm{continua en $I$} si lo es en todos sus puntos.
\end{definition}

Si $a$ no es un extremo de $I$, la condición dice exactamente que
\begin{keyformula}
  \lim_{x \to a} f(x) = f(a) ,
\end{keyformula}
\bymedium{%
  es decir: el límite existe y vale lo mismo que la función.
}{%
  es decir: el límite existe, $f$ está definida en $a$, y los dos valores
  coinciden. Si falla cualquiera de las tres cosas, $f$ no es continua en $a$.
}
\end{frame}

\begin{frame}
\didactatitle{Operaciones y discontinuidades}

\begin{remark}
Por el álgebra de límites, la suma, el producto y el cociente (donde el
denominador no se anula) de funciones continuas son continuos. En particular,
los polinomios y las funciones racionales son continuos allí donde están
definidos.
\end{remark}

\dpause

\begin{example}
\begin{itemize}
  \item $f(x) = \dfrac{x^2 - 1}{x - 1}$ para $x \neq 1$, con $f(1) = 0$. El
        límite en $1$ existe y vale $2$, pero no coincide con $f(1)$. La
        discontinuidad es \keyterm{evitable}: basta redefinir $f(1) = 2$.
  \item $H(x) = 0$ para $x < 0$ y $H(x) = 1$ para $x \geq 0$. No es continua
        en $0$: con $\varepsilon = 1/2$, cualquier intervalo alrededor de $0$
        contiene puntos $x < 0$, donde $|H(x) - H(0)| = 1$. Es un
        \keyterm{salto}, y ningún valor de $H(0)$ lo arregla.
\end{itemize}
\end{example}

\begin{commonmistake}
Confunden «continua» con «se dibuja sin levantar el lápiz», y no saben qué
hacer con $1/x$, que es continua en $(-\infty, 0)$ y en $(0, +\infty)$ aunque
su gráfica tenga dos ramas.
\end{commonmistake}
\end{frame}
